
\section{Related Work}

\subsection{Bottleneck Bandwidth and Round-Trip propagation Time (BBR)}

Bottleneck Bandwidth and Round-Trip propagation Time (BBR) was developed by Google researchers to address fundamental inefficiencies in traditional loss-based congestion control algorithms \cite{BBRORIGINAL}. 
The authors argue that modern hardware advancements have made the long-standing practice of using packet loss as a measure for congestion increasingly problematic \cite{BBRORIGINAL}. 
Specifically, in environments with large network buffers, loss-based algorithms like CUBIC tend to keep buffers full, leading to "bufferbloat" where delays are measured in seconds rather than milliseconds \cite{BBRORIGINAL}. 
Conversely, in long-distance networks, packet loss may occur due to link errors rather than actual congestion; however, loss-based algorithms misinterpret these signals, causing them to reduce throughput unnecessarily \cite{BBRORIGINAL}.

BBR operates through a state machine consisting of a Startup phase and two repeating steady-state phases known as ProbeBW and ProbeRTT. 
The authors define the optimal operating point as having a delivery rate equal to the bottleneck bandwidth ($BtlBw$) and a total amount of inflight data equal to the pipe's Bandwidth-Delay Product (BDP) \cite{BBRORIGINAL}. 
During the Startup phase, BBR rapidly discovers the bottleneck's limits by increasing its sending rate by a pacing gain of $2/\ln 2 \approx 2.885$ every round trip until the delivery rate plateaus \cite{BBRORIGINAL}.
While this ensures the bottleneck is discovered in $\log_2 BDP$ rounds, it may temporarily fill the queue with up to $2BDP$ of excess data before transitioning into a Drain phase to empty the buffer \cite{BBRORIGINAL, BBREVALUATION}.

In steady-state operation, the algorithm primarily cycles through the ProbeBW phase to measure and adapt to changes in available bandwidth. 
This phase uses an eight-step cycle where the algorithm typically increases its sending rate by a pacing gain of 1.25 for one RTT to probe for additional capacity, then immediately reduces it to a gain of 0.75 for one RTT to drain the queue created during the probe \cite{BBREVALUATION}. 
Finally, to refresh the minimum round-trip time estimate, BBR enters the ProbeRTT phase if a lower RTT has not been observed for 10 seconds \cite{BBREVALUATION}. 
During this phase, the algorithm abruptly limits inflight data to a minimal amount—typically 4 packets—for at least 200 ms to fully drain the bottleneck buffer and measure the true propagation delay of the path \cite{BBRORIGINAL, BBREVALUATION}.

\subsubsection{Active Queue Dept Compensation (AQDC) BBR}

BBR uses the ProbeBW phase to determine how much bandwidth it has.
During this time it increases its sending rate which will fill a queue, in case the available bandwidth has not increased.
The authors of \cite{AQDCBBR} see there a problem, because the queue may not be completly drained during the draining phase.
Furthermore, BBR is not fair to other flows, because of its filling up the queues during ProbeBW.
Another problem, that the authors saw is, that the algorithm uses the ProbeRTT phase to determine the latest RTT.
However, this occurs in fixed intervalls and therefore blindly, without considering how much build up in the queue is. 
\cite{AQDCBBR}

\subsubsection{Active Queue Debt Compensation (AQDC) BBR}

BBR utilizes the ProbeBW phase to estimate available bandwidth by periodically increasing its pacing rate. 
However, this process often injects bursts of traffic that exceed the path's instantaneous capacity, creating temporary queues \cite{AQDCBBR}. 
The authors of \cite{AQDCBBR} identify a critical limitation: these transient queues may not be completely drained during the subsequent ProbeDown phase, leading to a residual queue buildup that accumulates over successive cycles. 
Over time, this results in an upward drift of RTT measurements and increased tail latency, which is particularly detrimental to delay-sensitive applications like video streaming.

Furthermore, the authors contend that BBR's ProbeRTT mechanism is rigid and ``blind'' to queue history, as it resets at fixed intervals regardless of the actual backlog \cite{AQDCBBR}. 
To reconcile the conflict between fairness and low latency, AQDC-BBR explicitly models the ``queue debt'' ($B_{total}$) incurred during bandwidth probing. 
This debt is then compensated for by adjusting the target inflight data:

\begin{equation}
    Inflight_{target} = BDP - B_{total}
\end{equation}

By proportionally reducing the inflight data based on the tracked debt, AQDC-BBR drains accumulated queues and restores accurate minRTT tracking without the need for aggressive or heuristic rate cuts \cite{AQDCBBR}.

\subsubsection{BBR-Extended State (ES)}

\cite{BBRES}

\subsection{Performance-oriented Congestion Control (PCC)}

\cite{PCC}

\subsection{Copa}

\subsubsection{Copa-Dynamic (Copa-D)}

\cite{COPADYNAMIC}

\subsubsection{Copa+}

\cite{COPAPLUS}

\subsection{TCP's Third Eye}

\cite{TCPSTHIRDEYE}

\subsection{Survey on CCA}

\cite{SURVEYONCCA}

\subsection{Christian's Congestion Control Code - C4}

\cite{C4}
